\documentclass[11pt,a4paper]{article}
\usepackage[T1]{fontenc}
\usepackage[utf8]{inputenc}
\usepackage[margin=22mm]{geometry}
\usepackage{lmodern}
\usepackage{amsmath,amssymb}
\usepackage{enumitem}
\usepackage{xcolor}
\usepackage{microtype}
\usepackage{hyperref}
\definecolor{epflred}{HTML}{E3000B}
\hypersetup{colorlinks=true,urlcolor=epflred,linkcolor=epflred,
  pdftitle={ENG-654 Project 2: Geometry and Cuspidality in Orthogonal 3R Robots}}
\setlength{\parindent}{0pt}
\setlength{\parskip}{6pt}
\setlist[itemize]{leftmargin=*,itemsep=5pt,topsep=4pt}
\urlstyle{tt}

\begin{document}
{\small\textbf{EPFL \textbar\ ENG-654}\hfill Kinematics-Grounded Motion Planning for Robots}
\vspace{5mm}

{\color{epflred}\Large\bfseries Project 2}

{\LARGE\bfseries Geometry and Cuspidality in Orthogonal 3R Robots\par}

\textbf{Format:} Group of two students; simulation-based project.\\
\textbf{Prerequisites:} D--H models, position IK, Jacobians,
singularities, and configuration-space connectivity.

\section*{Motivation}
Two robots with the same number of joints can have very different global
inverse kinematics. A change in an axis offset can alter singularity curves,
the distribution of inverse-kinematic solutions, and the possibility of changing
posture without becoming singular. This project connects geometric design
choices to the joint-space and workspace maps needed for motion planning.

\section*{Task}
\textbf{Overall objectives.} Analyze a bounded family of orthogonal 3R position manipulators and explain
how its geometry controls singularities, two-/four-IK regions, and evidence of
cuspidal or noncuspidal behavior.

\begin{itemize}
  \item \textbf{Define and validate the family.} Use the idealized family
$a=(a_1,2,1.5)$, $d=(1,1,0)$, $\alpha=(\pi/2,\pi/2,0)$, with
$a_1\in\{0,0.5,1\}$. State units and derive forward position kinematics.
Check representative configurations using an independent implementation of the
stated transforms. The zero-offset endpoint is a limiting geometry of this
family, not the supplied custom 3R reference robot.

  \item \textbf{Construct reduced singularity plots.} Derive the position Jacobian
$J_p$ and check it by finite differences. Establish why $\det J_p$ is independent
of $q_1$. Plot its value and zero contours over $(q_2,q_3)\in[-\pi,\pi)^2$;
identify sampled connected regions without equating equal determinant signs
with connectivity.

  \item \textbf{Map the inverse distribution.} Map the singular curves to
$(\rho,z)$, where $\rho=\sqrt{x^2+y^2}$. Evaluate the supplied 3R IK method on
a workspace grid and distinguish zero, two, and four distinct real solutions.
Overlay the IK points of selected targets on the joint-space plot, including
targets on opposite sides of a critical curve.

  \item \textbf{Investigate posture connectivity.} For each geometry, select one target
with multiple regular IKs and examine whether two lie in the same computed
singularity-free component. Seek one nonsingular connecting path where supported
by the map. Explain what constitutes evidence of cuspidality and why failure
to find a path does not prove noncuspidality.

  \item \textbf{Compare geometry, not plot appearance.} Summarize changes in IK counts,
critical curves, and component membership. Refine the grids near a narrow region
or suspected cusp. Support any definitive classification with an analytical
argument and separate it from sampled observations; a numerical study alone
need not provide a definitive classification.
\end{itemize}

\newpage
\section*{Specifics and scope}
Use exactly three geometries and at least three diagnostic workspace targets
per geometry. Begin with approximately $150\times150$ joint and workspace grids
and refine one region of interest. Treat joint angles as periodic for the
unconstrained analysis. A complete symbolic classification of all orthogonal
3R geometries and a new quartic solver are not required.

Use the position Jacobian $J_p$ for the 3R task.
Confirm that the meridian reduction is valid for the selected model and
operational point. Count distinct real IKs modulo periodicity, then filter
joint limits separately. Identify two-/four-IK regions where present and
explain any absent count rather than forcing both types into every map. Show both determinant values and the zero contour,
and distinguish sampled connectivity from a proof. Refine the decisive
region instead of claiming completeness from a coarse grid.

\section*{Expected outcome}
Explain how D--H geometry changes global IK structure. Deliver three paired
joint/workspace maps, an IK-count comparison, and one detailed connectivity
example or a justified explanation of why the selected pair is separated.

Submit reproducible code, model parameters and test data, the required plots,
a concise 5--6-page report, and a short simulation demonstration. Explain
both successful cases and failures; conclusions must follow the numerical
and geometric evidence rather than an expected outcome.

\section*{Provided materials and implementation}
The custom 3R robot description (.urdf) and link meshes (.stl) will be provided
to the students, together with an analytical 3R IK implementation compatible
with the stated parameter family. The reference robot is orthogonal with
non-intersecting consecutive axes. It illustrates axis geometry; its dimensions
must not be substituted for those of the idealized family.

Use standard D--H transforms
$T_{i-1,i}=R_z(q_i)T_z(d_i)T_x(a_i)R_x(\alpha_i)$, with zero joint-angle
offsets for the idealized family. Implement the parameter study using these
transforms and use a plotting or robotics library to inspect frames, display
IKs and singularities, and animate the selected connection. Validate the
supplied IK results against your forward model.

\end{document}
